% TOP_PROBLEM: {"id": 3155, "title": "Domination in cubic graphs", "category_id": 3, "category": {"id": 3, "name": "graph_theory", "display_name": "Graph Theory", "description": "Problems involving graphs, networks, and their properties.", "slug": "graph-theory", "order_index": 3, "created_at": "2026-07-31T15:26:25.671Z"}, "status": "open", "problem_number": "OPG-37038", "set_id": 12}
% FIRST_POSTED: 2026-10-01
% ENTRY_KIND: statement_only
% UnsolvedMath ID 3155; source catalogue code OPG-37038
% !TeX program = lualatex
% Definition and sources only; this file is not an LLM solution attempt.
\documentclass[11pt,a4paper]{article}
\usepackage[margin=25mm]{geometry}
\usepackage{fontspec}
\setmainfont{lmroman10-regular.otf}[BoldFont=lmroman10-bold.otf,
 ItalicFont=lmroman10-italic.otf,BoldItalicFont=lmroman10-bolditalic.otf]
\usepackage{amsmath,amssymb,mathtools}
\usepackage{unicode-math}
\setmathfont{latinmodern-math.otf}
\AtBeginDocument{\let\setminus\smallsetminus}
\usepackage{xurl,hyperref}
\hypersetup{colorlinks=true,urlcolor=blue}
\setcounter{secnumdepth}{0}
\setlength{\parindent}{0pt}
\setlength{\parskip}{5pt}
\setlength{\emergencystretch}{3em}
\begin{document}
\section*{Domination in cubic graphs}\label{3155}
{\small UnsolvedMath ID 3155 · OPG-37038 · Graph Theory · OpenGarden}
\subsection{Definitions and mathematical statement}
Let $G=(V,E)$ be a finite simple undirected graph and write
$n=|V|$. For a vertex $v$, its closed neighbourhood is
$N[v]=\{v\}\cup\{u:uv\in E\}$. A set $D\subseteq V$ is dominating
if every vertex belongs to the closed neighbourhood of some vertex
of $D$, or equivalently
\[
 \bigcup_{v\in D}N[v]=V.
\]
The domination number $\gamma(G)$ is the minimum size of such a
set. Vertices in $D$ dominate themselves; there is no requirement
that they have another member of $D$ as a neighbour.

The graph is cubic if every vertex has degree three. It is
$3$-connected if $n\geq4$ and deleting any set of at most two
vertices leaves a connected graph. The proposed inequality is
\[
 \gamma(G)\leq\left\lceil\frac n3\right\rceil
 \qquad\text{for every $3$-connected cubic graph }G.
\]
Thus one asks for a set of at most the stated rounded-up size
meeting the closed neighbourhood of every vertex.

No independence or connectedness condition is imposed on $D$.
The graph need not be planar, bipartite, or have large girth.
Vertex connectivity is part of the hypothesis, so the problem
is not the assertion for arbitrary cubic graphs. The ceiling
has its ordinary integer meaning and matters at orders not
divisible by three. The quantification includes every admissible
finite order, rather than only an asymptotic estimate with an
unspecified additive error.
\subsection{Short English statement}
Can every 3-connected cubic graph be dominated by at most one third of its vertices, rounded up?
\subsection{Sources}
\begin{itemize}
\item[S1] ULAM.AI and the UnsolvedMath Contributors, supplied problems.json, record 3155 (OPG-37038), OpenGarden collection. Catalogue identity and source transcription; CC BY 4.0.
\url{https://huggingface.co/datasets/ulamai/UnsolvedMath}.
\item[S2] Open Problem Garden, Domination in cubic graphs. Original problem and discussion; the mathematical formulation below is expanded from the supplied source record.
\url{http://www.openproblemgarden.org/op/domination_in_cubic_graphs}.
\end{itemize}
\end{document}
